% 121-31(121쪽) — 두 곡선 y=x^2(x-4) 와 y=ax(x-4) 가 둘러싸는 두 도형. 스캔 화질이 낮아 TikZ 로 다시 그림.
% 원문에 있는 것만: x축 눈금 4 · 색칠한 두 도형 · 라벨 O, x, y, y=ax(x-4), y=x^2(x-4).
% a 는 원문 그림의 모양(포물선이 삼차곡선보다 얕고 두 도형이 둘의 교점에서 허리로 만난다)만 따라 그렸고 눈금을 넣지 않는다
% — a 가 답이므로 그림에서 값을 읽을 수 없어야 한다.
\begin{tikzpicture}[x=0.45cm, y=0.175cm, line width=0.5pt]
  % 색칠한 두 도형(원문)
  \path[fill=green!18]
    plot[domain=0:2, samples=60] (\x, {\x*\x*(\x-4)})
    -- plot[domain=2:0, samples=60] (\x, {2*\x*(\x-4)}) -- cycle;
  \path[fill=green!18]
    plot[domain=2:4, samples=60] (\x, {\x*\x*(\x-4)})
    -- plot[domain=4:2, samples=60] (\x, {2*\x*(\x-4)}) -- cycle;
  \draw[->, >=stealth, line width=0.7pt] (-1.1,0) -- (5.0,0) node[below=1pt] {$x$};
  \draw[->, >=stealth, line width=0.7pt] (0,-10.2) -- (0,7.2) node[left=1pt] {$y$};
  % 두 곡선
  \draw[line width=0.6pt, domain=-0.75:4.3, samples=150, smooth] plot (\x, {\x*\x*(\x-4)});
  \draw[line width=0.6pt, domain=-0.6:4.3, samples=120, smooth] plot (\x, {2*\x*(\x-4)});
  \node[above right=0pt and 1pt, font=\small] at (0,0) {$\pt{O}$};
  \node[below right=1pt and 0pt, font=\small] at (4,0) {$4$};
  \node[anchor=east, font=\small] at (4.6,8.3) {$y=ax(x-4)$};
  \node[anchor=north, font=\small] at (2.1,-10.6) {$y=x^2(x-4)$};
\end{tikzpicture}
